\scanpage{367}
\mspage{91}
admettant seulement $\mathrm{Loc}(7\mathrm{F})$, on sait qu'il existe une théorie des fibres $W$-homotopiques, mais on n'a \inserted{a priori} aucun moyen pratique pour ``calculer'' les\inserted{dites} fibres, faute de disposer d'un critère maniable \inserted{suffisant} pour assurer $p\in\mathrm{Fib}_W$. \deleted{\unsure{Je}}J'ai donc l'impression que l'axiome $\mathrm{Loc}(7\mathrm{F})$ ne va être d'un maigre usage. (Il est vrai que l'exemple de $\mathbb W_0$ devrait me rendre prudent ! Il me faudrait regarder le cas des $\mathbb W_{\unsure{n|m}}$, s'ils satisfont $\mathrm{Loc}(7\mathrm{F})$\ldots.)

À mon sens, l'axiome $\mathrm{Loc}(7\mathrm{D}_{W_{\unsure{\infty|\omega|w}}})$ est le \inserted{l'exact} ``minimum vital'' pour \deleted{\unsure{avoir}} être à l'aise dans le maniement des fibres $W$-homotopiques. Il est remarquable que ceci implique qu'on doive avoir $W\subset\mathbb W_\infty$, et que cela exclut du même coup \deleted{\uline{toutes}} les variantes considérées à ce jour de la notion du type d'homotopie !
% archive sheet number: 366
